\section{Introduction}
Multi-object tracking from an unmanned aerial vehicle is an observation problem
as much as an association problem. Density, target scale, occlusion, camera
motion, illumination, and blur change within a sequence, so a tracker receives
an observation stream whose quality and cost are both time-varying
\cite{zhu2022visdrone,du2018uavdt}. A fixed confidence floor can remove weak
but useful boxes in a sparse scene and retain clutter in a dense one; NMS can
merge nearby targets; and fixed resolution spends the same detector budget on
easy and difficult frames. These choices affect FP, FN, IDS, HOTA, MOTA, and
throughput together.

AC--MOT is a training-free runtime layer that computes inexpensive scene cues,
forms a Scene Complexity Index (SCI), and changes detector confidence, NMS, and
input resolution while keeping detector and tracker separate. Its scientific
value is the control interface: the detector operating point becomes an
auditable runtime variable rather than an immutable deployment setting.

The paper is organized around the optimization-driven evolution of that
interface. V1 was the first formal quality-oriented optimization;
V2 then made identity and deployment cost explicit through a multi-objective
Pareto study. Historical transfer and cross-pipeline studies exposed which
claims were portable and which were objective- or host-dependent. The modern
phase finally froze a stronger detector, a causal policy, a second tracker
host, runtime instrumentation, and attribution controls.

The contributions are bounded and evidence-oriented:
\begin{enumerate}
\item We present V1 as the first formal optimization-derived AC--MOT
      configuration and report its exact quality-oriented search and held-out
      uncertainty.
\item We document V2 as an identity-aware multi-objective extension with an
      explicit feasible Pareto front and an objective-dependent profile choice.
\item We report historical UAVDT transfer and separate U2MOT/VisDrone and
      SparseTrack/MOT17 cross-pipeline studies, including their negative or
      near-neutral evidence.
\item We specify the modern frozen AC--MOT v3 controller at implementation
      level: the
      crowd-selected SCI, surrounding state logic, thresholds, recovery probe,
      hysteresis, dwell, causal information flow, and action map.
\item We evaluate ByteTrack and a controlled OATrack host with aggregate metrics
      and paired 5000-resample sequence bootstrap intervals.
\item We measure end-to-end Tesla T4 latency and no-retuning UAVDT transfer
      for the two tracker hosts.
\end{enumerate}

The supported claim is precise. AC--MOT is an auditable training-free mechanism
for shaping the detector observation stream and computational operating point
of aerial MOT. The historical evidence supports distinct quality and
identity/runtime profiles; the frozen modern evidence supports positive changes
relative to the declared baseline and strong ByteTrack transfer, with a smaller
mixed effect on OATrack. It does not establish universal improvement, official
leaderboard equivalence, or that exact scene-conditioned timing is necessary
for the observed accuracy gains. Historical, cross-pipeline, and modern
protocols are therefore kept separate; extended historical and audit material
is provided in Online Resource 1.
